% Standalone mathematical section. Requires amsmath, amssymb, amsthm.
% Source-restoration theorem is supplied separately, with label compat:zero-restored.
% The ordinary scalar/iPRA bridge is inherited and must be included or cited.

\section{A mirror guard with one introduction rule}
\label{mir:section}

We combine a restriction of the formula grammar with a restriction of
equality transport. The orientation of the guard matters. Its initializer
is the \emph{left} side of the guarded equation; the iterator is the right
side. This permits the recovery of equation induction before transport
from a literal-zero source has been recovered.

\subsection{Complete syntax and primitive rules}

There is one numerical sort. Numerical terms are
\[
 t ::= x\mid 0\mid S(t)\mid I(n,a;u.h).
\]
Only the update body $h$ binds $u$. All other variables occurring in $h$
are fixed numerical parameters. Let
\[
 1:=S0,\qquad C(v,a,b):=I(v,a;u.b),
\]
where the displayed update variable is absent from $b$. These are
abbreviations, not additional primitive symbols.

Canonical formulas are generated by the following clauses.
\begin{enumerate}
\item Every numerical equation $s=t$ is a formula.
\item $D_{x,y}B$ is a formula if $B$ is a compound canonical formula.
\item $D_{x,y}B$ is a formula if its literal equation matrix has the form
\[
       B(x,y)\equiv a(x,y)=I(y,a(x,y);u.h(x,y,u)).
       \tag{Mirror matrix}\label{mir:mirror-matrix}
\]
\end{enumerate}
The bound coordinates $x,y$ are distinct. The finite terms $a,h$ may
also contain fixed parameters, and $a$ may depend on $y$. Bound variables
are renamed before substitution. Thus the occurrence of $a$ outside the
iterator and its initializer are literally the same term. The guard
condition is stable under capture-free substitution for free numerical
variables. The designated contradiction is $0=1$.

Contexts contain reusable clauses
\[
       \langle A;y.B(y)\rangle,
       \qquad y\notin\operatorname{FV}(A),
\]
with canonical $A,B$. An ordinary assumption $F$ abbreviates the constant
clause $\langle H;v.F\rangle$, where $v$ is fresh and
$H\equiv I(0,0;u.0)=0$. Its ordinary hypothesis use is the computation
of $H$ followed by the clause-use rule.

The two axiom schemas, in every context, are
\[
 I(0,a;u.h)=a,
 \qquad
 I(Sn,a;u.h)=h[I(n,a;u.h)/u].
 \tag{I0,IS}\label{mir:computation}
\]
There are precisely four other inference patterns:
\[
 \frac{\Gamma\vdash A}{\Gamma\vdash B(t)}
 \quad\text{if }\langle A;y.B(y)\rangle\in\Gamma
 \qquad(\mathrm{Use}),
\]
\[
 \frac{\Gamma\vdash s=t\qquad\Gamma\vdash F(s)}
      {\Gamma\vdash F(t)}
 \quad s\not\equiv0
 \qquad(\mathrm{Tr}^{-}),
 \tag{Restricted transport}\label{mir:transport}
\]
\[
 \frac{\Gamma\vdash D_{x,y}B(x,y)\qquad
       \Gamma,\langle B(b,0);y.B(b,y)\rangle\vdash A}
      {\Gamma\vdash A}\qquad(\mathrm{DE}),
\]
\[
 \frac{\Gamma,B(t,\eta_B(z))\vdash B(t,Sz)}
      {\Gamma\vdash D_{x,y}B(x,y)},
 \qquad
 \eta_B(z)=
 \begin{cases}
   z,&B\text{ is a literal equation},\\
   0,&B\text{ is compound}.
 \end{cases}
 \qquad(\mathrm{DI}).
 \tag{Root-determined introduction}\label{mir:introduction}
\]
In transport, $F$ is a canonical numerical formula template, with bound
variables fresh for the substituted terms. The restriction $s\not\equiv0$
is syntactic: variables and all $S$- or $I$-headed terms are permitted,
including those whose numerical value is zero. In DE, $b$ is fresh for
$\Gamma,A$ and the fixed parameters of $B$. In DI, $z$ is fresh for
$\Gamma,t$ and the fixed parameters of $B$. Displayed local clauses and
assumptions are discharged.

Call this calculus $\mathcal M$. There is no introduction-choice bit,
primitive reflexivity, successor axiom, coverage axiom, separation axiom,
or additional logical or induction rule. Weakening, ordinary-assumption
cut, and capture-free renaming are admissible by induction on derivations.
They preserve the syntactic restriction on transport sources. We do
\emph{not} assume arbitrary numerical substitution admissible at this
stage: replacing a variable source by literal zero need not preserve a
primitive transport instance.

\subsection{Equality operations available before zero-source transport}

\begin{lemma}[Safe equality operations]\label{mir:safe-equality}
Every term $r$ satisfies $\vdash r=r$. From $p=q$ with $p\not\equiv0$,
one may derive $q=p$ and $K(p)=K(q)$ for any fixed numerical template
$K$. Every computation equation may therefore be reversed as an
equation; its reversed direction is usable as a primitive transport
major only if its new source is not literally zero.
\end{lemma}
\begin{proof}
Let $b=I(0,r;u.0)$. The computation $b=r$, used both as transport
major and as minor in the template $v=r$, yields $r=r$. The source
$b$ has an iterator head. For symmetry, transport $p=p$ through
$v=p$ using $p=q$. For numerical congruence, transport
$K(p)=K(p)$ through $K(p)=K(v)$. Each source is the given nonzero
syntactic term $p$. Computation equations have an iterator-headed
left side, so the symmetry construction applies to them.
\end{proof}

\begin{lemma}[Joining nonliteral-zero terms at zero]\label{mir:zero-join}
If $p,q\not\equiv0$ and $\Gamma$ proves $p=0$ and $q=0$, then
$\Gamma\vdash p=q$, using computation and restricted transport only.
\end{lemma}
\begin{proof}
Put $z_0=I(0,0;u.0)$ and $b_p=I(0,p;u.0)$. From $b_p=b_p$,
transport with $p=0$ in the template $I(0,v;u.0)=b_p$ to obtain
$z_0=b_p$. Safe symmetry gives $b_p=z_0$. Transport using the
computation $b_p=p$ gives $p=z_0$. The identical construction gives
$q=z_0$. Safe symmetry at source $q$ gives $z_0=q$, and transport
at source $z_0$ in $p=z_0$ finishes. All sources used are among
$p,q,z_0,b_p,b_q$, hence are syntactically different from zero.
\end{proof}

\begin{lemma}[Expansion on the right]\label{mir:right-expansion}
Suppose $q\not\equiv0$ and $\Gamma\vdash q=r$. Under arbitrary
additional assumptions,
\[
                 s=q\quad\dashv\vdash\quad s=r.
\]
In particular, a constructor-headed computation can be expanded or
contracted on the \emph{right} of an equation without zero-source
transport, even when its computed value is literally zero.
\end{lemma}
\begin{proof}
Contraction uses $q=r$ in the template $s=v$, at the permitted source
$q$. For expansion, suppose $s=r$. If $r\not\equiv0$, reverse $q=r$
and transport at source $r$. If $r\equiv0$ and $s\not\equiv0$,
Lemma~\ref{mir:zero-join} joins $s=0$ and $q=0$. If both $r$ and $s$
are literally zero, safe symmetry of $q=0$ gives $0=q$ directly.
This is a finite case distinction on term syntax, not an object-level
case distinction or substitution into an existing derivation.
\end{proof}

\subsection{Guard endpoints, padding, and numerical cases}

For an equation $E\equiv(s=t)$ define its selector guard
\[
                  H_vE:=s=C(v,s,t).
                  \tag{Selector guard}\label{mir:selector-guard}
\]
This is the state-independent-update instance of
\eqref{mir:mirror-matrix} when used under a binder in $v$.

\begin{lemma}[Guard endpoints]\label{mir:guard-endpoints}
For every equation $E$, uniformly under retained contexts,
\[
                   \vdash H_0E,
                   \qquad H_{Sw}E\dashv\vdash E.
\]
Every canonical atomic matrix $B(x,y)$ also has the theorem $B(t,0)$
for every witness term $t$.
\end{lemma}
\begin{proof}
The base is the symmetry of the single zero computation
$C(0,s,t)=s$. For the successor endpoint, let $q=C(Sw,s,t)$.
The computation $q=t$ and Lemma~\ref{mir:right-expansion} give exactly
the displayed interderivability. For a general canonical atomic matrix,
its zero instance is
$a(t,0)=I(0,a(t,0);u.h(t,0,u))$, again the symmetry of one zero
computation. No congruence or equality-reflection theorem is assumed.
\end{proof}

\begin{lemma}[Case introduction at canonical matrices]
\label{mir:canonical-cases}
The case form of DI is admissible for every canonical matrix, including
literal equations.
\end{lemma}
\begin{proof}
For a compound matrix it is primitive. For an atomic canonical matrix
$B$, cut its theorem $B(t,0)$ into the given case premise
$\Gamma,B(t,0)\vdash B(t,Sz)$. Weaken the result by $B(t,z)$ and
apply the primitive progressive introduction. This establishes only
case introduction for matrices allowed by the canonical grammar.
\end{proof}

For an equation $E$ put
\[
 L(E):=D_{a,v}H_vE,
 \qquad
 P(A):=\begin{cases}L(A),&A\text{ is an equation},\\A,&A\text{ is compound}.
 \end{cases}
 \tag{Padding}\label{mir:padding-def}
\]
The new bound variables are fresh, and $a$ is unused. In particular,
$P(A)$ is always compound.

\begin{lemma}[Padding]\label{mir:padding}
$L(E)\dashv\vdash E$ and $P(A)\dashv\vdash A$ under arbitrary
retained contexts.
\end{lemma}
\begin{proof}
Under $E$, use guard introduction to prove $H_{Sz}E$ while ignoring
the local induction assumption $H_zE$. Atomic DI with witness zero
gives $L(E)$. Conversely, DE opens its clause
$\langle H_0E;v.H_vE\rangle$. The base is a theorem. Use at $1$
and successor guard removal give $E$. No numerical substitution into
an arbitrary premise derivation occurs.
\end{proof}

\begin{lemma}[Numerical cases]\label{mir:cases}
For any canonical formula $A(y)$, the rule
\[
 \frac{\Gamma\vdash A(0)\qquad\Gamma\vdash A(Sz)}
      {\Gamma\vdash A(t)}
 \tag{NC}\label{mir:nc}
\]
is admissible, provided $z$ is fresh for $\Gamma,t$ and the fixed
parameters of $A$.
\end{lemma}
\begin{proof}
If $A$ is compound, introduce $D_{u,y}A(y)$ by case DI with unused
$u$ and witness zero, ignoring its local base assumption. DE opens
$\langle A(0);y.A(y)\rangle$; the given base proof and Use at $t$
finish. If $A$ is an equation, pad the two premises, apply the compound
case rule to $L(A(y))$, and unpad its conclusion. The retained context
is unchanged throughout. In particular, this is not a rule permitting
case substitution in assumptions that contain the case variable.
\end{proof}

\subsection{A small implication sufficient for the bootstrap}

We need implication only between equations with a shared left side.
For fixed terms $s,a,b$, let
\[
 M_w(s;a,b):=s=C(w,a,b),\qquad
 \operatorname{Imp}_s(a,b):=D_{u,w}L(M_w(s;a,b)),
 \tag{Shared-left implication}\label{mir:shared-implication-def}
\]
with fresh controls and unused $u$. Lemma~\ref{mir:right-expansion}
applied to the zero and successor computations yields
\[
 M_0(s;a,b)\dashv\vdash(s=a),\qquad
 M_{Sz}(s;a,b)\dashv\vdash(s=b).
 \tag{Shared-left endpoints}\label{mir:shared-endpoints}
\]
The zero cases of the right endpoints are covered by zero-join. No
double-selector conversion that changes both sides of an equation is
being used.

\begin{lemma}[Rules for shared-left implication]\label{mir:shared-implication}
The following are admissible:
\[
 \frac{\Gamma,s=a\vdash s=b}
      {\Gamma\vdash\operatorname{Imp}_s(a,b)},
 \qquad
 \frac{\Gamma\vdash\operatorname{Imp}_s(a,b)\qquad\Gamma\vdash s=a}
      {\Gamma\vdash s=b}.
\]
\end{lemma}
\begin{proof}
For introduction, the outer matrix $L(M_w)$ is compound. Its case
assumption $L(M_0)$ unpads and contracts to $s=a$. Apply the supplied
derivation, expand $s=b$ to $M_{Sz}$ using
\eqref{mir:shared-endpoints}, and pad. Compound DI closes. For
elimination, DE opens the outer clause. Convert $s=a$ to its padded
zero antecedent, use the clause at $1$, unpad, and contract to $s=b$.
The supplied derivation is used at its original parameters; only fresh
control variables are introduced.
\end{proof}

\subsection{Full equation induction without substituting premise proofs}

\begin{theorem}[Contextual equation induction]\label{mir:equation-induction}
For every literal equation $E(y)$, $\mathcal M$ admits
\[
 \frac{\Gamma\vdash E(0)\qquad
       \Gamma,E(z)\vdash E(Sz)}
      {\Gamma\vdash E(t)},
 \tag{EI}\label{mir:ei}
\]
with the usual fresh-parameter condition on $z$. No transport from a
literal-zero source or numerical substitution into either supplied
premise derivation is used in its recovery.
\end{theorem}
\begin{proof}
Write $E(y)\equiv s(y)=r(y)$ and put
\[
 J(y):=s(y)=C(y,s(y),r(y)),
\]
\[
 \Phi(y):=
 \operatorname{Imp}_{s(y)}\bigl(C(y,s(y),r(y)),r(y)\bigr).
 \tag{Bootstrap implication}\label{mir:bootstrap-implication}
\]
Thus $\Phi(y)$ has the derived introduction and elimination rules of
$J(y)\to E(y)$. Its root is compound.

Under $\Gamma$, prove $\Phi(0)$ by shared-left implication introduction,
ignoring the assumption $J(0)$ and using the supplied base proof
$E(0)$. For a fresh $w$, prove $\Phi(Sw)$ by the same introduction:
its local hypothesis $J(Sw)$ contracts by the successor guard endpoint
to $E(Sw)$. These proofs retain only $\Gamma$, independent of the case
variable. NC therefore proves $\Phi(v)$ for every target term $v$.

In particular, under the new local assumption $J(z)$, the already
available $\Phi(z)$ gives $E(z)$. Apply the \emph{original} supplied
step derivation at this same eigenparameter $z$, obtaining $E(Sz)$.
The successor guard endpoint yields $J(Sz)$. Since $J(y)$ has the
canonical mirror shape, atomic progressive DI with an unused witness
coordinate gives $D_{u,y}J(y)$.

Eliminate this binder. Its clause has the theorem $J(0)$ as antecedent,
so Use at the desired target $t$ supplies $J(t)$. The independent NC
instance $\Phi(t)$ then gives $E(t)$ by shared-left implication
elimination. Close DE.

Every invocation of a guard or selector endpoint is a direct instance
of its uniformly constructed proof. In particular, the step derivation
has not been specialized at $0$, at $Sw$, or at any other new term.
This is the point that avoids circular reliance on unrestricted
numerical substitution in a source-restricted calculus.
\end{proof}

